% =====================================================================
%  Front matter -- Abstract (English, 200-350 words)
% =====================================================================
\chapter*{Abstract}
\addcontentsline{toc}{chapter}{Abstract}

Online services face traffic that rises and falls without warning, and the
tools that keep them fast are usually separate pieces. Classical load balancers
route requests by fixed rules that do not adapt to demand, cloud autoscalers add
capacity only after load has already grown, and monitoring and failure-handling
are bolted on afterwards. Putting these parts together gives operators many
dials but no single system that decides what to do. So latency spikes when
demand surges, and resources sit idle when it falls, which hurts most for small
and medium teams that cannot run heavy container-orchestration platforms.

This thesis presents SmartLoad, a lightweight middleware that brings demand
forecasting, anomaly detection, anomaly-aware request routing, and proactive
scaling into one decision loop, with a deterministic safety fallback on every action. SmartLoad
watches live telemetry from the backend pool, runs several decision engines in
parallel over a shared message bus, and applies a strict priority order: health
constraints first, learned recommendations second, and a classical rule as the
guaranteed fallback. Every automated decision is recorded and an operator can
override it at any time.

We built the system as a containerized stack and evaluated it on realistic
workloads. The forecast-driven scaling loop runs end to end: under rising load
the backend pool grew automatically up to its capacity ceiling, then shrank
again as demand fell, with the deterministic fallback ready on every step.
The deployed gradual-degradation detector closes a failure mode that point-based
rules miss, lifting its detection score from zero to \BManFone{} on slow drift.
SmartLoad's decisive advantage comes from anomaly-driven exclusion at equal capacity:
at the failure modes it targets, a failing backend that sheds errors but is
un-ejectable by passive health checks and a slow-but-healthy backend, it cut errors
roughly fivefold and tail latency by 8 to 250 times against static round-robin, robust
across three runs. Fine-grained routing is not where the gain lies: the deployed
heuristic ties round-robin under uniform overload, where an even split is provably
optimal, and a trained policy kept dormant in shadow mode proves round-robin-equivalent,
reported as an honest null result rather than a learned-routing win.

We conclude that an integrated, safety-first decision loop is a practical
foundation for adaptive load management: the closed loop works end to end, the
measured results are reported clearly, and richer learned behaviour on highly
varied workloads is a natural next step.

\clearpage
